\documentclass{article}
\usepackage{amsmath, amssymb, amsthm}
\usepackage{hyperref}
\usepackage{geometry}
\geometry{margin=1in}

\title{Erd\H{o}s Problem \#669}
\date{Last updated: 2025-08-31}
\author{Source: erdosproblems.com}

\begin{document}
\maketitle

\noindent\textbf{Status:} OPEN \quad \textbf{No prize}

\noindent\textbf{Tags:} geometry

\medskip
\noindent\textbf{Note:} orchard problems

\medskip

\noindent\textbf{Problem Statement:}

Let $F_k(n)$ be minimal such that for any $n$ points in $\mathbb{R}^2$ there exist at most $F_k(n)$ many distinct lines passing through at least $k$ of the points, and $f_k(n)$ similarly but with lines passing through exactly $k$ points. Estimate $f_k(n)$ and $F_k(n)$ - in particular, determine $\lim F_k(n)/n^2$ and $\lim f_k(n)/n^2$.




Trivially $f_k(n)\leq F_k(n)$ and $f_2(n)=F_2(n)=\binom{n}{2}$. The problem with $k=3$ is the classical 'Orchard problem' of Sylvester. Burr, Gr\"{u}nbaum, and Sloane \cite{BGS74} have proved that\[f_3(n)=\frac{n^2}{6}-O(n)\]and\[F_3(n)=\frac{n^2}{6}-O(n).\]There is a trivial upper bound of $F_k(n) \leq \binom{n}{2}/\binom{k}{2}$, and hence\[\lim F_k(n)/n^2 \leq \frac{1}{k(k-1)}.\]See also [101] .




References

[BGS74] Burr, Stefan A. and Gr\"{u}nbaum, Branko and Sloane, N. J. A., The orchard problem . Geometriae Dedicata (1974), 397-424.

\noindent\textbf{Related OEIS sequences:} A003035, A006065, A008997, possible


\bigskip
\noindent\small{Source: \url{https://www.erdosproblems.com/669}}
\end{document}
